\section*{Chapitre \ref{ch:b1:fluid-envelopes-structure} --- La structure de l'atmosphère et de l'océan}
% values: calc/ch26.py

\begin{solution}{exo:b1:fluid-envelopes-structure:1}
(a)~$H = R^*T/(M_0 g_0) = 8314.32 \times 250/(28.9644 \times 9.80665) =
\qty{7.32}{km}$. (b)~$e^{10/7.32} = 3.92$~: la pression chute à environ le quart.
(c)~$H$ est proportionnelle à la température~: à \qty{216.65}{K} elle vaut
\qty{6.34}{km}, contre \qty{8.43}{km} à \qty{288.15}{K}. L'air froid est plus dense à
une pression donnée, donc la pression chute plus vite avec la hauteur.
\end{solution}

\begin{solution}{exo:b1:fluid-envelopes-structure:2}
Avec $\Gamma_d = \qty{9.76}{K/km}$~: \qtyrange{0}{0.3}{km}, $11.0 > 9.76$,
instable~; \qtyrange{0.3}{1.2}{km}, $9.7$, neutre aux incertitudes des
données près~; \qtyrange{1.2}{1.5}{km}, une inversion~; \qtyrange{1.5}{8}{km}, stable~;
au-dessus de \qty{8}{km}, isotherme, stable. L'inversion est la plus stable, puisque
$\Gamma_d - \Gamma = 13.8$ y est le plus grand.
\end{solution}

\begin{solution}{exo:b1:fluid-envelopes-structure:3}
(a)~$\sigma = 1026.95 - 1000 = \qty{26.95}{kg.m^{-3}}$. (b)~B, la plus légère, se trouve
au-dessus de A. (c)~Toutes les masses volumiques de l'eau de mer commencent par les mêmes chiffres, 10~; $\sigma$ ne garde
que les chiffres qui changent d'une eau à l'autre.
\end{solution}

\begin{solution}{exo:b1:fluid-envelopes-structure:4}
(a)~$S = 0.25 \times 35.40 + 0.75 \times 34.30 = 34.575$~; $\theta = 0.25 \times
15.0 + 0.75 \times 4.0 = \qty{6.75}{\celsius}$. (b)~Sur le segment PQ, au quart
du chemin de Q à P (\cref{prop:b1:fluid-envelopes-structure:mixing}).
(c)~Les mélanges de P, Q et R remplissent le triangle PQR. En écrivant
$(34.60, 8.0) = a\,\mathrm P + b\,\mathrm Q + c\,\mathrm R$ avec $a + b + c = 1$
on trouve $a = 0.345$, $b = 0.788$, $c = -0.133$~: une proportion négative de R, donc le
point est hors du triangle et aucun mélange des trois ne peut le produire.
\end{solution}

\begin{solution}{exo:b1:fluid-envelopes-structure:5}
(a)~$\theta = 253.15 \times 2^{2/7} = \qty{308.6}{K}$. (b)~À \qty{1000}{hPa},
$\theta = T = \qty{298.15}{K}$. La \omterm{def:b1:fluid-envelopes-structure:potential-temperature}{température potentielle} croît vers le haut, de
298.2 à \qty{308.6}{K}~: la couche est stable pour de l'air sec. (c)~Il arriverait à
sa \omterm{def:b1:fluid-envelopes-structure:potential-temperature}{température potentielle}, \qty{308.6}{K}, soit \qty{35.4}{\celsius}, dix
degrés plus chaud que l'air déjà présent.
\end{solution}

\begin{solution}{exo:b1:fluid-envelopes-structure:6}
(a)~$T = 288.15 - 6.5 \times 5 = \qty{255.65}{K}$ et $p = \num{101325} \times
(255.65/288.15)^{5.256} = \qty{54020}{Pa}$, \qty{540.2}{hPa}. (b)~$H =
\qty{8.43}{km}$ et $p = \num{101325}\,e^{-5/8.43} = \qty{56010}{Pa}$. La pression
réelle (standard) est plus basse~: l'air y est plus froid que \qty{288.15}{K}, donc
plus dense, et la pression chute plus vite avec la hauteur. (c)~$p = \num{22632}\,
e^{-9/6.34} = \qty{5475}{Pa}$, la valeur du
\cref{tab:b1:fluid-envelopes-structure:layers}.
\end{solution}

\begin{solution}{exo:b1:fluid-envelopes-structure:7}
(a)~$\beta\,\Delta S = \alpha\,\Delta T$ donne $\Delta S = 1.67/7.61 = 0.22$.
(b)~$\Delta\rho = 1026.95 \times (-1.67\times10^{-4} \times 2 + 7.61\times10^{-4}
\times 0.3) = \qty{-0.11}{kg.m^{-3}}$~: l'eau devient plus légère.
(c)~$\Delta\rho = 1028.11 \times (-0.53\times10^{-4} \times 2 + 7.85\times10^{-4}
\times 0.3) = \qty{+0.13}{kg.m^{-3}}$~: les mêmes variations la rendent maintenant plus dense, parce que
le réchauffement ne change presque pas la masse volumique de l'eau proche du point de congélation.
\end{solution}

\begin{solution}{exo:b1:fluid-envelopes-structure:8}
(a)~$N^2 = (9.81/260) \times (9.76 - 4)\times10^{-3} = \qty{2.17e-4}{s^{-2}}$,
$N = \qty{0.0147}{s^{-1}}$, période $2\pi/N = \qty{426}{s}$, environ sept
minutes. (b)~$N^2 = (9.81/270) \times 14.76\times10^{-3} = \qty{5.36e-4}{s^{-2}}$,
$N = \qty{0.0232}{s^{-1}}$, période \qty{271}{s}. (c)~$N$ tend vers zéro et la
période croît sans borne~: la force de rappel s'évanouit dans une couche neutre.
\end{solution}

\begin{solution}{exo:b1:fluid-envelopes-structure:9}
(a)~Elle est plus dense, $\sigma_\theta = 27.20$ contre 26.30~: l'AAIW est plus douce
mais bien plus froide. (b)~$\Delta\sigma = 1026.95 \times [-1.67\times10^{-4} \times
(4.31 - 15.43) + 7.61\times10^{-4} \times (34.30 - 35.54)] = 1.907 - 0.969 =
0.94$, contre $0.90$ exactement. (c)~La température~: le refroidissement ajoute
\qty{1.91}{kg.m^{-3}}, le dessalement retire \qty{0.97}{kg.m^{-3}}.
\end{solution}

\begin{solution}{exo:b1:fluid-envelopes-structure:10}
(a)~D'après le \cref{thm:b1:fluid-envelopes-structure:stability}, $(g/\theta)\,\mathrm
d\theta/\mathrm dz = (g/T)(\Gamma_d - \Gamma)$~; $\mathrm d\theta/\mathrm dz = 0$
donne $\Gamma = \Gamma_d$. (b)~$T = T_s - \Gamma_d z$, une couche de gradient constant
$L = \Gamma_d$~: d'après la \cref{prop:b1:fluid-envelopes-structure:layer-pressure}, $p =
p_s (T/T_s)^{gM/(R\Gamma_d)}$, et $gM/(R\Gamma_d) = gMc_p/(Rg) = c_pM/R$
($= 3.5$ pour l'air sec). (c)~$p$ s'annule là où $T = 0$, en $z = T_s/\Gamma_d =
300/9.76 = \qty{30.7}{km}$. Les couches réelles de \omterm{def:b1:fluid-envelopes-structure:potential-temperature}{température potentielle} uniforme font
au plus quelques kilomètres d'épaisseur~: elles sont coiffées d'air stable, comme la troposphère l'est
par la stratosphère.
\end{solution}

\begin{solution}{exo:b1:fluid-envelopes-structure:11}
(a)~$S = (33.787 + 36.050)/2 = 34.919$, $\theta = \qty{8.0}{\celsius}$.
(b)~Avec une équation d'état linéaire, $\sigma$ est une fonction affine de $S$ et
de $\theta$, donc le $\sigma$ du mélange est la moyenne des deux, 27.00.
(c)~$\alpha$ vaut \qty{0.75e-4}{K^{-1}} à \qty{2}{\celsius} et
\qty{2.07e-4}{K^{-1}} à \qty{14}{\celsius}, tandis que $\beta$ change peu
(\num{7.80e-4} et \num{7.53e-4}). La masse volumique chute de plus en plus vite quand la
température monte~: c'est une fonction concave de la température. Le long de la ligne
de mélange, le mélange est à la température moyenne, où la masse volumique se situe au-dessus de la
corde joignant les deux parents, c'est-à-dire au-dessus de 27.00. (d)~L'excès vient de la
courbure, un terme du second ordre~: il croît comme le carré de la différence de température
des parents.
\end{solution}

\begin{solution}{exo:b1:fluid-envelopes-structure:12}
(a)~La différence de pression est $5604 - 5500 = \qty{104}{dbar}$~: la parcelle se refroidit
de $0.129 \times 0.104 = \qty{0.013}{\celsius}$, jusqu'à \qty{0.718}{\celsius}.
(b)~L'eau alentour à \qty{5400}{m} est à \qty{0.730}{\celsius}, plus chaude~: la
parcelle, de salinité presque identique, est plus dense que son environnement et retombe.
La couche est stable, bien que la température mesurée croisse
vers le bas~; sa \omterm{def:b1:fluid-envelopes-structure:potential-temperature}{température potentielle} décroît vers le bas. (c)~$\theta \approx
0.7313 - 0.086 \times 5.604 = \qty{0.249}{\celsius}$, à un millième de
degré de la valeur exacte.
\end{solution}

\begin{solution}{pb:b1:fluid-envelopes-structure:1}
\begin{enumerate}
  \item La température est presque uniforme jusqu'à \qty{50}{m}, chute vite jusqu'à
        \qty{300}{m}, plus lentement jusqu'à \qty{1000}{m}, et est presque uniforme en dessous.
        La salinité décroît de 35.80 en surface à un minimum de 34.32 à
        \qty{800}{m}, puis croît jusqu'à 34.90 à \qty{3000}{m}.
  \item $22.40 - 0.5 = \qty{21.90}{\celsius}$ est atteint entre \qty{50}{m}
        (22.10) et \qty{75}{m} (19.20)~: $50 + 25 \times 0.20/2.90 = \qty{52}{m}$.
  \item $(22.10 - 13.00)/250 = \qty{0.0364}{\celsius/m}$, \qty{3.6}{\celsius}
        par \qty{100}{m}. La \omterm{def:b1:fluid-envelopes-structure:ocean-layers}{thermocline} s'étend de la base de la \omterm{def:b1:fluid-envelopes-structure:ocean-layers}{couche de
        mélange}, vers \qty{50}{m}, à environ \qty{800}{m}~; elle est la plus marquée entre
        \qty{50}{m} et \qty{75}{m} (\qty{11.6}{\celsius} par \qty{100}{m}).
  \item $\sigma = 1026.95 \times [1 - 1.67\times10^{-4} \times 12.4 +
        7.61\times10^{-4} \times 0.80] - 1000 = 25.45$, \qty{0.74}{kg.m^{-3}} trop
        haut. L'eau de surface est loin de l'état de référence, et $\alpha$
        est bien plus grand à \qty{22}{\celsius} qu'à \qty{10}{\celsius}~: la
        forme linéaire sous-estime l'allègement dû à la chaleur.
  \item Le minimum de salinité à \qty{800}{m}, 34.32 à \qty{4.8}{\celsius}, est
        l'Eau Intermédiaire Antarctique, formée sous les pluvieux vents d'ouest de l'océan
        Austral. Le maximum à \qty{3000}{m}, 34.90 à
        \qty{2.5}{\celsius}, est l'Eau Profonde Nord-Atlantique, formée dans les
        mers du Groenland et du Labrador.
  \item $(2.50 - 2.26)/3.042 = \qty{0.079}{\celsius}$ par \qty{1000}{dbar}.
  \item $\sigma_\theta$ croît avec la profondeur~: la colonne est stable. La
        \omterm{def:b1:fluid-envelopes-structure:ocean-layers}{pycnocline} est là où $\sigma_\theta$ croît le plus vite, dans les quelques
        centaines de mètres supérieurs, avec la \omterm{def:b1:fluid-envelopes-structure:ocean-layers}{thermocline}~: \num{0.82} entre
        \qty{200}{m} et \qty{800}{m}, seulement \num{0.70} sur les
        \qty{2200}{m} en dessous.
  \item $S = (35.45 + 34.32)/2 = 34.885$~; $\theta = (14.90 + 4.80)/2 =
        \qty{9.85}{\celsius}$.
  \item C~: 26.461~; I~: 27.310~; mélange~: 26.886, exactement la moyenne des deux
        (26.886), puisque la forme linéaire est affine.
  \item La moyenne des parents est $(26.342 + 27.159)/2 = 26.751$~; le mélange,
        à 26.888, est plus dense de \qty{0.14}{kg.m^{-3}}. C'est le cabbeling.
  \item Avec la forme linéaire, la descente donne 26.61 à \qty{300}{m}, 26.88 à
        \qty{500}{m} et 27.31 à \qty{800}{m}~: la masse volumique du mélange se situe près de
        \qty{500}{m}. L'eau en ce lieu, $(34.82, \qty{9.6}{\celsius})$, est très
        proche du mélange $(34.885, \qty{9.85}{\celsius})$~: la descente entre
        l'Eau Centrale et l'AAIW est elle-même proche d'une ligne de mélange entre
        elles.
  \item $S = 0.25 \times 35.45 + 0.75 \times 34.32 = 34.60$, $\theta = 0.25
        \times 14.90 + 0.75 \times 4.80 = \qty{7.33}{\celsius}$. L'interpolation de
        la descente entre \qty{500}{m} et \qty{800}{m} donne
        $(34.57, \qty{7.2}{\celsius})$ vers \qty{650}{m}~: presque dessus.
  \item En \si{K/km}~: A--B 10.0~; B--C 9.7~; C--D $-5.5$~; D--E 6.29~; E--F 6.50~;
        F--G 6.57~; G--H 0.33~; H--I $-0.50$~; I--J $-0.83$.
  \item A--B instable~; B--C neutre~; C--D une inversion~; D--E, E--F, F--G stables~;
        G--H presque isotherme, stable~; H--I et I--J des inversions. La tropopause
        est en G, \qty{12.5}{km} et \qty{-59.0}{\celsius}, là où la température
        cesse de décroître.
  \item La température moyenne est \qty{214.65}{K}~: $H = 287 \times 214.65/9.81 =
        \qty{6.28}{km}$, et $147\,e^{-4/6.28} = \qty{77.7}{hPa}$, cohérent
        avec les 78 mesurés.
  \item L'air en A, \qty{22.0}{\celsius}, est plus frais que la mer~: la chaleur passe
        de la mer à l'air, qui est chauffé par en dessous. La couche A--B, avec
        un gradient vertical supérieur à $\Gamma_d$, en est le résultat.
  \item En B~: $22.0 - 9.76 \times 0.5 = \qty{17.1}{\celsius}$, contre 17.0~; en
        C~: \qty{7.4}{\celsius}, contre 7.3~; en D~: \qty{3.5}{\celsius}, contre
        9.5. La parcelle est légèrement plus chaude que l'air jusqu'en C, et bien plus froide
        en D.
  \item À la base de l'inversion, près de \qty{1.5}{km}~: en entrant dans l'inversion,
        la parcelle devient plus froide, donc plus dense, que l'air alentour, ce qui
        arrête son ascension.
  \item Les parcelles libérées par la couche de surface instable la traversent en montant et la mélangent
        jusqu'à ce que sa \omterm{def:b1:fluid-envelopes-structure:potential-temperature}{température potentielle} soit uniforme, ce qui rend son gradient vertical égal à
        $\Gamma_d$~; l'inversion au-dessus arrête le mélange.
  \item En K~: A 294.15~; B 293.99~; C 294.08~; D 300.51~; E 312.65~; F 329.29~; G
        346.28~; H 369.50~; I 446.99~; J 600.91.
  \item $\theta$ décroît de A à B (instable), est presque constant de B à
        C (neutre) et croît au-dessus (stable), le plus vite de C à D et au-dessus de
        G~: la même classification qu'à la question~14.
  \item E--F~: $(329.29 - 312.65)/4 = \qty{4.2}{K/km}$~; H--I~: $(446.99 -
        369.50)/4 = \qty{19.4}{K/km}$, près de cinq fois plus.
  \item Avec $\theta = \qty{408.2}{K}$ en moyenne~: $N^2 = (9.81/408.2) \times
        0.01937 = \qty{4.66e-4}{s^{-2}}$, $N = \qty{0.0216}{s^{-1}}$, une période
        de \qty{291}{s}, environ cinq minutes.
  \item $N^2 = (9.81/1027) \times (26.35 - 24.79)/150 = \qty{9.9e-5}{s^{-2}}$,
        $N = \qty{0.0100}{s^{-1}}$, une période de \qty{630}{s}. La basse
        stratosphère est la plus stable, avec un $N^2$ près de cinq fois plus grand
        que dans la \omterm{def:b1:fluid-envelopes-structure:ocean-layers}{thermocline}.
  \item En G, $\theta = 214.15 \times (1000/186)^{2/7} = \qty{346}{K}$~: la
        \omterm{def:b1:fluid-envelopes-structure:potential-temperature}{température potentielle} à la tropopause est d'environ \qty{346}{K}.
\end{enumerate}
\end{solution}
